\chapter{Introducción}

La criptografía simétrica protege información mediante algoritmos que comparten una clave secreta. Su seguridad práctica depende, entre otros factores, de que un adversario no pueda recuperar la clave ni distinguir el comportamiento del cifrador con recursos razonables. El criptoanálisis diferencial estudia cómo ciertas diferencias entre entradas se propagan hasta las salidas y permite construir ataques estadísticos contra versiones reducidas o configuraciones vulnerables de algunos cifradores. La disponibilidad de algoritmos cuánticos obliga a revisar el coste de estos ataques bajo modelos de acceso distintos a los clásicos \citep{biham1991,kaplan2017,zhou2019}.

Este Trabajo estudia una aplicación cuántica del análisis diferencial basada en la búsqueda de candidatos de subclave que maximizan una puntuación de conteo. El problema se formula a partir de pares de textos cifrados con diferencias relacionadas; para cada candidato se cuenta cuántos pares satisfacen una condición de filtrado derivada de la estructura del cifrador. Una búsqueda de mínimo o máximo puede reducir el número de consultas al conjunto de candidatos, aunque el coste de calcular cada puntuación y de implementar su oráculo reversible debe contabilizarse por separado \citep{durr1999,zhou2019}.

\section{Motivación y justificación}

La relevancia del problema reside en comprender qué parte de un ataque diferencial puede acelerarse mediante computación cuántica y bajo qué supuestos. Una cota de consultas cuadrática no implica por sí sola una reducción equivalente del tiempo total: el algoritmo debe acceder a los datos de forma coherente, evaluar reversiblemente la condición de filtrado y deshacer los cálculos auxiliares. La literatura diferencia, además, entre adversarios que ejecutan cómputo cuántico sobre datos clásicos y adversarios capaces de realizar consultas cuánticas al cifrador; ambos modelos no son intercambiables \citep{kaplan2017}.

\section{Planteamiento del problema}

Se considera un ataque de última ronda en el que se recuperan \(b\) bits de subclave mediante una colección de \(N\) pares de textos. Para cada candidato \(k\in\{0,1\}^{b}\), sea \(f_i(k)\in\{0,1\}\) el indicador de que el par \(i\), tras el descifrado parcial correspondiente, satisface el filtro diferencial. La puntuación es
\begin{equation}
 A(k)=\sum_{i=1}^{N} f_i(k), \qquad k\in\{0,1\}^{b}.
 \label{eq:score}
\end{equation}
El objetivo es identificar un candidato con puntuación máxima o suficientemente alta. La pregunta de investigación es: ¿bajo qué supuestos puede una búsqueda cuántica de máximo reducir el coste de explorar los \(2^b\) candidatos, y qué coste total resulta al incluir la evaluación de \(A(k)\), los datos y los requisitos de acceso coherente?

\section{Alcance y estructura del Trabajo}

El TFE estudia la recuperación de una subclave parcial de última ronda mediante criptoanálisis diferencial y la aplicación de una búsqueda cuántica de máximo a la función de puntuación de los candidatos. El estudio comprende la formulación del ataque, la definición del modelo de acceso a los datos y al cifrador, el análisis de los requisitos de un oráculo reversible y la comparación de los costes clásico y cuántico en términos de datos, consultas, operaciones y memoria. La evaluación computacional se circunscribe a una instancia reducida y reproducible, con parámetros explícitos, para contrastar la función de puntuación y las estimaciones de coste.

Quedan fuera del alcance la recuperación de la clave completa, la evaluación exhaustiva de la seguridad de cifradores comerciales y la generalización de los resultados a configuraciones no estudiadas. Tampoco se abordan la ejecución en hardware cuántico, el ruido, la corrección de errores ni una estimación de recursos físicos. Las conclusiones sobre una posible ventaja cuántica se limitan al modelo de acceso y a los costes analizados; una reducción de consultas no se interpreta por sí sola como una ventaja práctica de tiempo de ejecución.

El Capítulo 2 revisa el contexto y los antecedentes; el Capítulo 3 fija los objetivos y la metodología; el Capítulo 4 define el modelo del problema, los requisitos y los supuestos; el Capítulo 5 presenta el diseño de la solución y el análisis de costes; el Capítulo 6 aborda la evaluación computacional; y el Capítulo 7 expone las conclusiones, las limitaciones y el trabajo futuro.
